\chapter{Ultrasound}

\section*{What Is This Test?}
Ultrasound uses high-frequency sound waves to create real-time images of internal organs. It's safe (no radiation), shows movement (blood flow, heart beating), and is excellent for soft tissues. Common uses: abdomen, pelvis, thyroid, blood vessels, pregnancy.

\section*{Alternative Names}
US, sonography, ultrasonography, ultrasound scan, sonogram, echography


\section*{Principle}
Ultrasound uses high-frequency sound waves (2-18 MHz) generated by piezoelectric crystals in a transducer. The waves propagate through tissue and reflect at interfaces between tissues of different acoustic impedance. The reflected waves (echoes) are received by the same crystals and converted to a real-time grayscale image.

\section*{History}
Medical ultrasound was pioneered by Ian Donald in the 1950s in Glasgow. The first obstetric ultrasound was performed in 1958, and real-time imaging became available in the 1970s, revolutionizing prenatal care.

\section*{Why This Test Is Ordered}
Your doctor orders an ultrasound when they need to see inside your body without radiation. Ultrasound is used for many purposes: evaluating abdominal organs, checking for gallstones, examining the thyroid, guiding needle procedures, assessing blood flow, and monitoring pregnancy. It answers the question: what does the structure look like, and is there fluid, a mass, or abnormal blood flow?

\section*{Is This a Primary or Secondary Test}
A primary test for many conditions. Often the first imaging test because it is safe, fast, and provides real-time information.

\section*{How This Test Is Performed}
A technician applies a gel to your skin and moves a handheld probe over the area being examined. The probe sends sound waves into your body and creates images based on the echoes that come back. You may feel light pressure from the probe but there is no pain. The test typically takes thirty to sixty minutes depending on what is being examined.

\section*{What Preparation Is Needed}
Preparation depends on what is being examined:
\begin{itemize}
\item Abdominal ultrasound: You may need to fast for six to eight hours beforehand so your gallbladder and pancreas can be seen clearly.
\item Pelvic ultrasound: You may need a full bladder to provide a window for viewing your reproductive organs.
\item Thyroid and vascular ultrasound: No preparation is needed.
\item Pregnancy ultrasound: No preparation is needed.
\end{itemize}

\section*{Contraindications and Precautions}
Diagnostic ultrasound has no absolute contraindications. It uses no ionizing radiation and is safe for pregnant women and children. For endoscopic ultrasound (EUS), precautions related to sedation apply.

\section*{Risks}
Diagnostic ultrasound carries no known biological risks. It is a highly safe, radiation-free imaging modality. For endoscopic ultrasound (EUS), there is a small risk of bleeding or perforation, especially if a biopsy is performed.

\section*{What Happens After the Test}
There is no recovery time. You can wipe off the gel and resume normal activities immediately. Results may be available within minutes if the ultrasound is performed by your doctor, or within a few hours if interpreted by a radiologist.

\section*{Understanding Your Results}

\section*{Normal Results}
A normal result shows no abnormalities in the imaged organs. Normal size, shape, echotexture, and blood flow patterns.

\section*{Follow-up Tests}
If the ultrasound is abnormal, your doctor may order a CT scan or MRI for more detailed evaluation, a Doppler ultrasound to assess blood flow, a contrast-enhanced ultrasound, or a biopsy guided by ultrasound.
